\documentclass[../../../include/open-logic-section]{subfiles}

\begin{document}

\olfileid[fr]{sth}{choice}{tarskiscott}
\olsection{L'astuce de Tarski--Scott}

Dans la
\olref[cardinals][cardsasords]{defcardinalasordinal},
nous avons défini les cardinaux comme des
ordinaux. Pour cela, nous avons supposé
l'axiome du bon ordre. Nous l'avons fait
pour la seule raison que c'est l'usage courant.

Avant d'aborder les questions philosophiques
entourant le bon ordre, il est donc important
de préciser que nous \emph{pouvons} nous écarter
de cet usage et élaborer une théorie des
cardinaux \emph{sans} supposer l'axiome du
bon ordre. Nous pouvons conserver les
définitions de $\cardeq{A}{B}$,
$\cardle{A}{B}$ et $\cardless{A}{B}$
données dans le \olref[sfr][siz][]{chap}.
Il nous faudra seulement une nouvelle
notion de \emph{cardinal}.

Une idée naïve serait de définir la cardinalité
de $A$ par :
\[
	\Setabs{x}{\cardeq{A}{x}}.
\]
Vous pouvez comparer cette idée à la définition
frégéenne de $\# x Fx$, esquissée tout à la
fin de la \olref[cardinals][hp]{sec}.
Pour les raisons que nous y avons évoquées,
cette définition échoue. Tout singleton est
équipotent à $\{\emptyset\}$. Or de nouveaux
singletons apparaissent à chaque étape
successeur de la hiérarchie : il suffit
de considérer le singleton de l'étape
précédente. Ainsi, si $A$ est un singleton,
$\Setabs{x}{\cardeq{A}{x}}$ ne peut être
un ensemble, car cette collection n'aurait
pas de rang.\footnote{Note de l'édition
française : l'original affirme sans restriction
que cette collection n'existe pas. Pour
$A=\emptyset$, elle vaut pourtant
$\{\emptyset\}$ ; l'exemple des singletons
suffit à invalider la définition générale.}

Pour contourner ce problème, nous utilisons
une astuce due à Tarski et Scott :\footnote{
Rappel : toutes les formules peuvent comporter
des paramètres, sauf indication contraire.}

\begin{defn}[Tarski--Scott]
Pour toute formule $\phi(x)$, soit
$[ x : \phi(x)] $ l'ensemble de tous les $x$
de plus petit rang possible tels que $\phi(x)$
(ou $\emptyset$, s'il n'existe aucun $x$
satisfaisant $\phi$).
\end{defn}

Vérifions que cette définition est légitime.
Dans $\ZF$,
le \olref[spine][foundation]{zfentailsregularity}
garantit que $\setrank{x}$ existe pour
chaque $x$. S'il existe des objets satisfaisant
$\phi$, nous pouvons prendre pour $\alpha$
le plus petit rang tel que
$(\exists x\subseteq V_\alpha)\phi(x)$,
c'est-à-dire tel que
$(\forall \beta \in \alpha)
(\forall x \subseteq V_\beta)\lnot \phi(x)$.
Nous pouvons alors définir
$[x : \phi(x)]$ par séparation comme
$\Setabs{x \in V_{\alpha+1}}{\phi(x)}$.

L'astuce de Tarski--Scott étant justifiée,
nous pouvons l'utiliser pour définir
une notion de cardinalité :

\begin{defn}
La \emph{cardinalité \textsc{ts}} de $A$ est
$\text{tsc}(A) = [x : \cardeq{A}{x}]$.
\end{defn}

La définition d'un cardinal \textsc{ts}
n'utilise pas l'axiome du bon ordre.
Mais, même sans cet axiome, nous pouvons
montrer que les \emph{cardinaux \textsc{ts}}
se comportent à peu près comme les
\emph{cardinaux} définis dans la
\olref[cardinals][cardsasords]{defcardinalasordinal}.
Par exemple, si nous reformulons
le \olref[cardinals][cardsasords]{lem:CardinalsBehaveRight}
et
le \olref[card-arithmetic][opps]{lem:SizePowerset2Exp}
en termes de cardinaux \textsc{ts},
leurs démonstrations restent valables
dans $\ZF$, sans supposer l'axiome du
bon ordre.

Profitons-en pour signaler que l'astuce
de Tarski--Scott permet aussi de développer
une théorie des ordinaux. Si
$\tuple{A, <}$ est un bon ordre, posons
$\text{tso}(A, <) =
[\tuple{X, R} :
\ordeq{\tuple{A, <}}{\tuple{X, R}}]$.
Pour en savoir plus sur cette façon de
traiter les cardinaux et les ordinaux,
voir \citet[chs.~9--12]{Potter2004}.

\end{document}
